\documentclass[11pt,a4paper]{article}

\usepackage{kotex}
\usepackage{geometry}
\usepackage{graphicx}
\usepackage{booktabs}
\usepackage{array}
\usepackage{hyperref}
\usepackage{caption}
\usepackage{longtable}
\usepackage{amsmath}
\usepackage{setspace}

\geometry{margin=1in}
\setstretch{1.15}

\title{10기가비트 이더넷에서 433 ns 지연시간을 달성한 FPGA 기반 고빈도 거래 시스템\\
\large 원문 한국어 번역본}
\author{Yi-Chieh Kao, Hung-An Chen, Hsi-Pin Ma\\
Department of Electrical Engineering, National Tsing Hua University}
\date{}

\begin{document}
\maketitle

\begin{abstract}
고빈도 거래 시스템은 시장 데이터 피드에 반응하여 수익을 내기 위해 극도로 낮은 지연시간을 요구한다. 지연시간을 줄이기 위해 본 시스템은 25 ns의 낮은 지연시간을 갖는 10기가비트 이더넷 물리 트랜시버, 맞춤형 네트워크 스택 파싱 및 패키징, 부분적인 금융 프로토콜 디코딩 및 인코딩, 호가창 처리, 맞춤형 거래 전략을 사용하여 구현되었다. 거래 시스템의 하드웨어 테스트와 기능 검증은 Yuanta Futures가 제공한 대만 선물 거래 환경을 사용하였다. 내부 시장 패킷 분석부터 주문 패킷 발생까지의 지연시간은 약 433 ns이다.
\end{abstract}

\noindent\textbf{키워드:} 고빈도 거래, 하드웨어 가속, 선물

\section{서론}

\subsection{배경}
지난 20년 동안 기술적 발전과 규제 변화는 금융시장이 작동하는 방식을 크게 바꾸어 놓았다. 시장은 완전히 사람 중심으로 운영되던 시스템에서 주로 컴퓨터화된 거래 시스템으로 변화하였다.

고빈도 거래(HFT)는 극도로 낮은 지연시간을 가진 정교한 컴퓨터 프로그램을 사용하여 매우 짧은 시간 안에 막대한 수의 거래 체결을 수행하고, 매우 빈번한 거래와 작은 스프레드를 통해 수익을 축적하는 방식을 의미한다. 따라서 수많은 주문이 금융시장으로 전송되며, 왕복 시간은 마이크로초 단위 이내로 추정된다. 그러나 주어진 기회에서 수익을 얻을 수 있는 것은 거래소에 먼저 요청을 보낸 소수의 거래자뿐이다. 이것이 지연시간이 중요한 주된 이유이다.

HFT가 가져오는 큰 수익 때문에 점점 더 많은 거래자가 HFT 영역에 참여하고 이러한 기회를 포착하려고 한다. 그 결과 거래 속도 경쟁은 점점 더 치열해지고 있다.

\subsection{주요 특징}
\begin{itemize}
    \item 네트워크 스택 디코더는 FPGA 위에서 수신 프레임을 효율적인 파이프라인 구조로 디코딩한다. 이를 통해 이더넷 프레임이 NIC에서 수신된 뒤 운영체제에 의해 디코딩되는 일반적인 소프트웨어 기반 시스템의 높은 지연시간을 줄인다.
    \item 대만 선물 시장을 위한 맞춤형 금융 프로토콜 디코더를 구현하였다. 원래의 파싱 과정을 단순화하고, 거래 동작에 영향을 주는 메시지에만 집중한다.
    \item 본 논문에서 제안한 호가창은 단일 상품에 대해 매수 및 매도 양쪽의 상위 5개 가격과 수량, 그리고 파생 가격을 처리할 수 있다. 범용 호가창과 비교하여 호가창 처리 지연시간을 최소화한다. 또한 여러 상품을 처리하기 위해 호가창 모듈을 추가하기만 하면 되는 확장 가능한 하드웨어 구조를 제시한다.
    \item 맞춤형 TCP 오프로딩 엔진을 구현하여 TCP 연결과 연결 해제의 기본 기능을 제공한다.
    \item DMA 서브시스템을 통해 서로 다른 연결 파라미터에 기반한 수동 주문 입력 및 TCP 동작 기능을 제공한다.
\end{itemize}

\section{하드웨어 아키텍처}
제안된 10기가비트 이더넷 고빈도 거래 시스템의 구조는 원문 그림 1에 제시되어 있다.

첫 번째 부분은 네트워크 계층 모듈이다. 여기에는 TCP/IP 참조 모델에서의 네트워크 접근 계층, 인터넷 계층, 전송 계층 기능이 포함된다. 이 모듈은 네트워크로부터 패킷을 수신하거나 네트워크로 패킷을 송신하며, TCP/IP 또는 UDP/IP 패킷을 금융 프로토콜 디코더로 오프로딩한다.

두 번째 부분은 금융 프로토콜 디코더와 인코더이다. 금융 프로토콜 디코더는 TAIFEX 메시지 프로토콜(TMP)과 대만 틱 바이 틱 데이터 프로토콜이라는 두 종류의 금융 프로토콜 파싱을 수행한다. 네트워크 계층에서 전달된 페이로드를 금융 시장 정보로 처리하고, 특정 상품의 현재 상태를 호가창에서 갱신한다. 또한 거래소 서브시스템 로그인에 필요한 연결 절차를 트리거하기 위해 금융 프로토콜 인코더로 정보를 보낸다. 금융 프로토콜 인코더는 거래자가 직접 만든 수동 주문 요청 또는 자동 맞춤형 거래 전략의 출력을 TCP 패킷으로 변환한다.

그 다음 호가창은 특정 상품의 상위 5개 매수/매도 가격을 처리하는 데 사용된다. 맞춤형 거래 로직은 호가창의 정보를 바탕으로 스프레드를 얻기 위해 주문을 제출한다. 마지막으로 DMA 서브시스템 또는 맞춤형 거래 로직에서 발생한 주문 정보를 저장하기 위해 주문 FIFO가 시스템에 추가된다.

시간적으로 중요하지 않은 작업에는 DMA 서브시스템이 사용되어 호스트 CPU와 FPGA 사이의 데이터 전송 기능을 제공한다. 따라서 감독, 보고와 같은 작업은 호스트 CPU에서 수행될 수 있다. 이를 통해 제안 시스템은 PCIe의 예측 불가능성을 피할 수 있다.

\end{document}
